\section{A variational equation for all Gamma coefficients of derivative order at most two}
\label{var:section}

This section answers the next-sector question in
\cite[research questions 4--5]{PTI}. The result is independent of the
tail determinant: it is an exact identity in a free differential
polynomial algebra. After substitution of digamma derivatives, it
provides coefficients in identities among Gamma, polygamma, and
height-one zeta integrals.

\subsection{The inherited normal form and the new finite-jet principle}

Let
\[
 \mathscr R=\Q[u_0,u_1,u_2,\ldots],\qquad
 Du_r=u_{r+1},\qquad \wt(u_r)=r+1.
\]
A monomial is \emph{functional} if it is a power of $u_0$, including
$1$, or if its largest positive derivative index occurs at least twice.
Write $\mathscr F$ for their rational linear span.
The classical functional-monomial decomposition
\cite[Proposition 4.10]{Guo2014}, used in \cite{DPL,PTI}, is
\begin{equation}\label{var:normal}
 P=\NN(P)+D\HH(P),\qquad \NN(P)\in\mathscr F,
 \qquad [1]\HH(P)=0.
\end{equation}
It is unique. For completeness, if a monomial $M$ has unique largest
index $r\ge1$, replace $u_r$ by $u_{r-1}$ to form $T$, and let $a$
be the multiplicity of $u_{r-1}$ in $T$. Then
\[
 M=\frac1aDT+\left(M-\frac1aDT\right).
\]
Every remaining monomial has smaller largest derivative index.
This proves existence by termination. If $H$ is nonconstant and has
largest derivative index $r$, the nonzero term
$u_{r+1}\partial H/\partial u_r$ in $DH$ has unique largest index;
it cannot belong to $\mathscr F$. This proves uniqueness. Both maps
preserve factor degree, while $D$ raises weight by one. The quotient
is a vector-space quotient: $\NN$ is linear but is not multiplicative.
For example, $\NN(u_1)=0$ while $\NN(u_1^2)=u_1^2$.

Define the Euler variational derivative
\begin{equation}\label{var:Euler}
 \mathscr E(P)=\sum_{r\ge0}(-D)^r\frac{\partial P}{\partial u_r}.
\end{equation}
The sum is finite for each polynomial. The commutator identity
$\partial_{u_r}D=D\partial_{u_r}+\partial_{u_{r-1}}$, with the
last term zero for $r=0$, proves by cancellation that
$\mathscr E(DH)=0$.

For $d\ge0$ let $\sigma_d$ set all $u_r$, $r>d$, equal to zero,
and let
\[
 \delta_d=\sum_{r=0}^{d-1}u_{r+1}\partial_{u_r}
\]
act on $\Q[u_0,\ldots,u_d]$.

\begin{lemma}[Finite-jet variational principle]\label{var:finitejet}
For every polynomial $P$,
\begin{equation}\label{var:finitejet-eq}
 \sigma_d\mathscr E(P)=
 \sum_{r=0}^d(-\delta_d)^r
      \partial_{u_r}\bigl(\sigma_d\NN(P)\bigr).
\end{equation}
Thus the variational derivative determines a closed equation for a
truncated normal form even though normal reduction itself must precede
truncation.
\end{lemma}
\begin{proof}
Equation \eqref{var:normal} and $\mathscr E D=0$ permit replacing
$P$ by $\NN(P)$ on the left. Any normal monomial removed by
$\sigma_d$ has its highest index greater than $d$ at least twice.
Any partial derivative in \eqref{var:Euler} removes at most one such
factor. Every subsequent application of $D$ preserves at least one
factor of index greater than $d$. Its specialization is therefore zero.
For the retained terms, the ideal generated by $u_{d+1},u_{d+2},\ldots$
is $D$-stable, so specialization of $D^r$ is $\delta_d^r$.
These observations prove the identity term by term.
\end{proof}

\subsection{The inverse-Gamma hierarchy}

Set
\begin{equation}\label{var:hierarchy}
 \Acal_0=1,\qquad \Acal_{m+1}=u_0\Acal_m-D\Acal_m.
\end{equation}
The Bell exponential is
\begin{equation}\label{var:Bell}
 \Acal(t)=\sum_{m\ge0}\Acal_m\frac{t^m}{m!}
 =\exp\left(\sum_{r\ge0}
         \frac{(-1)^ru_rt^{r+1}}{(r+1)!}\right).
\end{equation}
Under $u_r=\psi^{(r)}(x)$, these become
\[
 \Acal_m(x)=(-1)^m\Gamma(x)\frac{d^m}{dx^m}\frac1{\Gamma(x)},
 \qquad \Acal(t,x)=\frac{\Gamma(x)}{\Gamma(x-t)}.
\]
All operations below are formal in $t$; no interchange with an
infinite integral is needed for the coefficient theorem.

Translation $u_0\mapsto u_0+c$ commutes with $D$ and preserves
$\mathscr F$, hence commutes with $\NN$. It follows that
$\NN(\Acal(t))$ has an Appell factor $e^{u_0t}$.
There is also the centered identity
\begin{equation}\label{var:centered}
 \NN(\Acal(t))=
 \NN\exp\left(\sum_{j\ge0}
          \frac{u_{2j}t^{2j+1}}{2^{2j}(2j+1)!}\right).
\end{equation}
Indeed the exponential on the right before normal reduction is
$e^{tD/2}\Acal(t)$, and every coefficient of
$(e^{tD/2}-1)\Acal(t)$ is a total derivative.
Consequently every retained monomial has factor degree congruent to
its weight modulo two.

There are uniquely determined formal series $f_b(y)$ such that
\begin{equation}\label{var:fdefinition}
 \sigma_2\NN(\Acal(t))
 =e^{u_0t} f(y,z),\qquad
 f(y,z)=\sum_{b\ge0}z^bf_b(y),\qquad
 y=\frac{u_1t^2}{2},\quad z=u_2t^3.
\end{equation}
Each $f_b$ is even in $y$. The functional-monomial convention excludes
one copy of $u_2$, so $f_1=0$.

\begin{theorem}[Complete three-variable Gamma-sector equation]
\label{var:master}
Put $L=2y+(z/2)\partial_y$. Then
\begin{equation}\label{var:PDE}
 \boxed{\displaystyle
 f-\frac12L f_y+L^2f_z
 =e^{z/24}\int_0^1
          \cosh(yv)e^{zv^2/8}\dd v.}
\end{equation}
The equation, the normal-form conditions $f_1=0$ and absence of a
linear $y$ term, and the normalization $f_0(0)=1$ uniquely determine
all the functions in \eqref{var:fdefinition}.
More explicitly, let
\begin{equation}\label{var:Rn}
 R_n(y)=\frac1{n!24^n}\int_0^1(1+3v^2)^n\cosh(yv)\dd v.
\end{equation}
Then
\begin{align}
 f_0(y)&=\frac12\left(
       \cosh y+\frac{\sinh y}{y}-y\Shi(y)\right),
       & f_1(y)&=0,\label{var:f0}\\
 4(n+1)y^2f_{n+1}(y)
 &=R_n(y)-(n+1)f_n(y)-(2n-1)y f_n'(y)
           -\frac{n-2}{4}f_{n-1}''(y),
       && n\ge1.\label{var:recurrence}
\end{align}
Here $\Shi(y)=\int_0^y\sinh v\dd v/v$. Every apparent singularity
at $y=0$ in these formulas is removable.
\end{theorem}
\begin{proof}
Differentiating \eqref{var:Bell} with respect to a variable $u_r$
and applying the Euler operator gives
\begin{equation}\label{var:EA}
 \mathscr E\Acal(t)
 =\sum_{r\ge0}\frac{t^{r+1}}{(r+1)!}D^r\Acal(t)
 =\int_0^t e^{sD}\Acal(t)\dd s.
\end{equation}
After setting $u_r=0$ for $r\ge3$, a formal shift replaces $u_0$
by $u_0+s u_1+s^2u_2/2$ and $u_1$ by $u_1+s u_2$.
The exponent of $e^{sD}\Acal(t)$ becomes
\[
 tu_0+(st-t^2/2)u_1+
       (ts^2/2-t^2s/2+t^3/6)u_2.
\]
Substitute $s=t(1+v)/2$. The exponent is
$tu_0+yv+z/24+zv^2/8$. Integrating over $-1\le v\le1$
and using its even part yields the right side of
\eqref{var:PDE}, multiplied by $te^{tu_0}$.

On $e^{tu_0}g(y,z)$ the restricted derivation is
$e^{tu_0}t^{-1}Lg$. Therefore Lemma~\ref{var:finitejet} gives
\[
 \sigma_2\mathscr E\NN\Acal(t)
 =te^{tu_0}\{f-\tfrac12Lf_y+L^2f_z\}.
\]
This proves \eqref{var:PDE}. Expanding
\[
 L^2=4y^2+z+2zy\partial_y+\frac{z^2}{4}\partial_y^2
\]
and comparing $z^n$ gives \eqref{var:recurrence}. At $z^0$ the
condition $f_1=0$ gives
$f_0-yf_0'=\sinh y/y$. Its even solution with constant one is
\eqref{var:f0}; the only homogeneous freedom is a multiple of $y$,
which the normal form excludes.

For $n\ge1$, \eqref{var:recurrence} determines $f_{n+1}$ uniquely
because multiplication by $y^2$ is injective on formal series.
Divisibility of its numerator by $y^2$ follows from the existence of
the normal form already proved in \eqref{var:normal}. This is not an
assumption about arbitrary initial functions. The coefficient theorem
below also proves that each $f_b$ is entire in $y$.
\end{proof}

\subsection{Rational formulas for every coefficient}

For even $k\ge0$, write
\[
 r_b(k)=k![y^k]f_b(y).
\]
Then $r_0(k)=-1/((k-1)(k+1))$ and $r_1(k)=0$.
Taking coefficients in \eqref{var:recurrence} gives the explicit
rational recurrence
\begin{equation}\label{var:rrecurrence}
\begin{split}
 r_{n+1}(k)=\frac1{4(n+1)(k+1)(k+2)}\bigg\{
 &\frac1{n!24^n}\sum_{j=0}^n
              \binom nj\frac{3^j}{k+2j+3}\\
 &-\bigl(n+1+(2n-1)(k+2)\bigr)r_n(k+2)
 -\frac{n-2}{4}r_{n-1}(k+4)\bigg\}.
\end{split}
\end{equation}
This computes a whole coefficient family with rational operations,
without expanding all partitions at its total weight.

\begin{theorem}[Polynomial numerators and fixed odd denominators]
\label{var:rational}
For every $b\ge2$ there is a rational polynomial $p_b(X)$ of degree
exactly $b-2$ such that
\begin{equation}\label{var:rationalform}
 \boxed{\displaystyle
 r_b(k)=\frac{p_b(k)}{\prod_{j=0}^{b}(k+2j+1)},\qquad
 [X^{b-2}]p_b(X)=\frac1{4b!6^{b-1}}.}
\end{equation}
These formulas hold for all even $k\ge0$. In particular every
$f_b(y)$ is entire. If $a,k\ge0$ and $b\ge2$, then
\begin{equation}\label{var:allcoefficients}
 [u_0^au_1^ku_2^b]\NN(\Acal_{a+2k+3b})
 =\begin{cases}
 \displaystyle\frac{(a+2k+3b)!}{a!2^kk!}
       \frac{p_b(k)}{\prod_{j=0}^b(k+2j+1)},& k\text{ even},\\[5pt]
 0,&k\text{ odd}.
 \end{cases}
\end{equation}
A coefficient with one copy of $u_2$ is zero.
\end{theorem}
\begin{proof}
The case $b=2$ follows from \eqref{var:rrecurrence}:
$p_2=1/48$. For a polynomial recurrence, set
\[
 \Delta_n(X)=\prod_{h=0}^n(X+2h+3),\qquad
 S_n(X)=\frac1{n!24^n}\sum_{j=0}^n
       \binom nj3^j\frac{\Delta_n(X)}{X+2j+3}.
\]
Each quotient in this expression is a polynomial. For $n\ge2$,
\eqref{var:rrecurrence} becomes
\begin{equation}\label{var:precurrence}
\begin{split}
 p_{n+1}(X)=\frac1{4(n+1)(X+2)}\bigg\{
 S_n(X)&-\bigl(n+1+(2n-1)(X+2)\bigr)p_n(X+2)\\
 &-\frac{n-2}{4}(X+3)p_{n-1}(X+4)\bigg\},
\end{split}
\end{equation}
where the last term is zero for $n=2$.
The numerator vanishes at $X=-2$: after division by
$\Delta_n(-2)\ne0$, that assertion is exactly the constant-term
identity in \eqref{var:recurrence}, whose left side contains $y^2$.
Thus \eqref{var:precurrence} is a polynomial recurrence, not merely
a rational-function definition.

The degree of $S_n$ is $n$, and its leading coefficient is
$4^n/(n!24^n)=1/(n!6^n)$. The two subtracted terms have degree at
most $n-1$ and $n-2$. Division by $4(n+1)(X+2)$ proves that
$p_{n+1}$ has degree $n-1$ and the stated leading coefficient.
Induction proves \eqref{var:rationalform}.
For fixed $b$, $r_b(k)=O(k^{-3})$ as $k\to\infty$, and the
factorial in the series for $f_b$ proves entire convergence.
Finally substitute $y=u_1t^2/2$, $z=u_2t^3$ in
\eqref{var:fdefinition} and expand $e^{u_0t}$.
The factor-degree parity already proved gives the zero odd coefficients.
\end{proof}

The first polynomial numerators are
\begin{align*}
 p_2(X)&=\frac1{48},&
 p_3(X)&=\frac{2X+5}{1728},\\
 p_4(X)&=\frac{X^2+7X+9}{20736},&
 p_5(X)&=\frac{2X^3+27X^2+100X+90}{1244160},\\
 p_6(X)&=\frac{2X^4+44X^3+316X^2+832X+603}{44789760}.
\end{align*}
The theorem supplies every subsequent numerator by the same finite
recurrence. It does not require that all its displayed possible poles
be actual poles; some factors can cancel at higher $b$.

\subsection{Closed square and cube sectors}

\begin{corollary}[All mixed tetragamma-square and cube coefficients]
\label{var:squarecube}
The two first new entire functions are
\begin{align}
 f_2(y)&=\frac{(y^2+3)\sinh y-3y\cosh y}{48y^5}
       =\frac1{384}\int_0^1(1-v^2)^2\cosh(yv)\dd v,
       \label{var:f2closed}\\
 f_3(y)&=\frac{(2y^4+60y^2+135)\sinh y
                -15y(y^2+9)\cosh y}{1728y^7}.
       \label{var:f3closed}
\end{align}
In particular, for all $a,s\ge0$,
\begin{align}
 [u_0^au_1^{2s}u_2^2]\NN(\Acal_{a+4s+6})
 &=\frac{(a+4s+6)!}
 {48a!2^{2s}(2s)!(2s+1)(2s+3)(2s+5)},
 \label{var:squarecoeff}\\
 [u_0^au_1^{2s}u_2^3]\NN(\Acal_{a+4s+9})
 &=\frac{(4s+5)(a+4s+9)!}
 {1728a!2^{2s}(2s)!(2s+1)(2s+3)(2s+5)(2s+7)}.
 \label{var:cubecoeff}
\end{align}
The removable values are $f_2(0)=1/720$ and $f_3(0)=1/36288$.
\end{corollary}
\begin{proof}
The numerator formulas for $p_2,p_3$ give their power series directly.
For the square sector use
\[
 \int_0^1v^k(1-v^2)^2\dd v
 =\frac8{(k+1)(k+3)(k+5)}.
\]
For the cube sector use
\[
 f_3(y)=\frac{2y\partial_y+5}{82944}
                \int_0^1(1-v^2)^3\cosh(yv)\dd v.
\]
Repeated elementary integration gives \eqref{var:f2closed} and
\eqref{var:f3closed}. Coefficient extraction in
\eqref{var:allcoefficients} proves the final formulas.
\end{proof}

For example, the coefficients of $u_2^3$ in $\NN(\Acal_9)$,
of $u_1^2u_2^2$ in $\NN(\Acal_{10})$, and of $u_2^4$ in
$\NN(\Acal_{12})$ are respectively $10,90,$ and $220$.
At weight $14$, the coefficient of $u_1^4u_2^2$ is $15015$.
These are consequences of the all-order formulas, not the evidence
from which their validity is inferred.

\subsection{What the sector theorem does and does not evaluate}

The specialization $\sigma_2$ is performed after normal reduction.
Substituting $u_r=\psi^{(r)}(x)$ into a full moment identity does not
make the omitted derivatives vanish. For instance
\[
 \NN(\Acal_8)
 =u_0^8-70u_0^4u_1^2+28u_0^2u_2^2-7u_1^4-u_3^2.
\]
The last term remains in the digamma identity. The new theorem solves
the coefficient problem for a complete sector and supplies a
finite-jet variational method for higher sectors; it does not claim
that the sector alone evaluates an arbitrary pure digamma moment.
The next section states the exact transfer with its full boundary
functional.
